\chapter{Συμπεράσματα}

Η παρούσα εργασία υλοποίησε και αξιολόγησε ένα πειραματικό πλαίσιο ανίχνευσης ανωμαλιών σε ακολουθίες καταγραφών \en{HDFS}, χρησιμοποιώντας το \en{bert-base-uncased} με προσαρμογή \en{LoRA}. Το πρόβλημα διατυπώθηκε ως δυαδική ταξινόμηση ακολουθιών, με στόχο τη διάκριση φυσιολογικής και ανώμαλης συμπεριφοράς. Η αξιολόγηση έγινε σε δύο βασικές συνθήκες: κεντρική εκπαίδευση, όπου το μοντέλο είχε πρόσβαση σε ολόκληρο το σύνολο εκπαίδευσης, και ομοσπονδιακή εκπαίδευση, όπου το σύνολο εκπαίδευσης κατανεμήθηκε σε 2, 4 και 8 \en{clients}. Στα ομοσπονδιακά πειράματα εξετάστηκαν οι στρατηγικές \en{FedAvg} και \en{FedAdam}, καθώς και δύο τιμές της παραμέτρου \(\alpha\) του \en{Dirichlet partitioning}, \(\alpha=0.8\) και \(\alpha=100\).

\section{Κύρια Ευρήματα}

Η κεντρική εκπαίδευση έδωσε την υψηλότερη συνολική επίδοση. Στο σύνολο ελέγχου πέτυχε \en{accuracy} \(0.9663\), \en{precision} \(0.9717\), \en{recall} \(0.9308\) και \en{F1-score} \(0.9494\). Το αποτέλεσμα αυτό αποτελεί το ισχυρό σημείο αναφοράς της εργασίας, καθώς προκύπτει από εκπαίδευση πάνω στο ενιαίο σύνολο δεδομένων, χωρίς τον κατακερματισμό και τη συνάθροιση ενημερώσεων που εισάγει η ομοσπονδιακή μάθηση.

Η υψηλή τιμή \en{precision} δείχνει ότι το κεντρικό μοντέλο παράγει σχετικά λίγες λανθασμένες προβλέψεις ανωμαλίας, ενώ η υψηλή τιμή \en{recall} δείχνει ότι εντοπίζει μεγάλο μέρος των πραγματικών ανωμαλιών. Η μικρή απόσταση μεταξύ \en{precision} και \en{recall} οδηγεί σε υψηλό \en{F1-score} και δείχνει ισορροπημένη συμπεριφορά. Αυτό είναι ιδιαίτερα σημαντικό για την ανίχνευση ανωμαλιών σε αρχεία καταγραφής, όπου τόσο τα ψευδώς θετικά όσο και τα ψευδώς αρνητικά αποτελέσματα έχουν πρακτικό κόστος.

Στα πειράματα με 2 \en{clients}, οι τελικές τιμές \en{F1-score} στο σύνολο ελέγχου κυμάνθηκαν από \(0.8129\) έως \(0.8201\). Η καλύτερη ρύθμιση ήταν η \en{FedAvg} με \(\alpha=100\), η οποία πέτυχε \en{accuracy} \(0.8765\), \en{precision} \(0.8245\), \en{recall} \(0.8160\) και \en{F1-score} \(0.8201\). Οι διαφορές ανάμεσα στις τέσσερις εκτελέσεις ήταν μικρές, γεγονός που δείχνει ότι, για μικρό πλήθος \en{clients}, η επιλογή στρατηγικής συνάθροισης και η αλλαγή της παραμέτρου \(\alpha\) δεν οδήγησαν σε μεγάλη διαφοροποίηση της τελικής απόδοσης.

Στα πειράματα με 4 \en{clients}, οι τελικές τιμές \en{F1-score} στο σύνολο ελέγχου κυμάνθηκαν από \(0.8138\) έως \(0.8245\). Η καλύτερη ρύθμιση ήταν η \en{FedAdam} με \(\alpha=100\), με \en{accuracy} \(0.8776\), \en{precision} \(0.8236\), \en{recall} \(0.8253\) και \en{F1-score} \(0.8245\). Η συγκεκριμένη εκτέλεση αποτέλεσε την καλύτερη συνολική ομοσπονδιακή επίδοση της εργασίας. Σε αυτή τη ρύθμιση, η προσαρμοστική ενημέρωση του διακομιστή μέσω \en{FedAdam} φαίνεται να αξιοποιεί καλύτερα τις τοπικές ενημερώσεις σε σχέση με τον απλό σταθμισμένο μέσο όρο.

Στα πειράματα με 8 \en{clients}, οι τελικές τιμές \en{F1-score} στο σύνολο ελέγχου κυμάνθηκαν από \(0.8011\) έως \(0.8222\). Η καλύτερη ρύθμιση ήταν η \en{FedAdam} με \(\alpha=100\), με \en{accuracy} \(0.8732\), \en{precision} \(0.8150\), \en{recall} \(0.8303\) και \en{F1-score} \(0.8222\). Παρότι η καλύτερη τιμή \en{F1-score} ήταν ελαφρώς χαμηλότερη από την αντίστοιχη των 4 \en{clients}, η συγκεκριμένη εκτέλεση πέτυχε την υψηλότερη τιμή \en{recall} μεταξύ όλων των ομοσπονδιακών πειραμάτων. Αυτό δείχνει ότι η αύξηση του πλήθους των \en{clients} μπορεί να ευνοήσει τον εντοπισμό περισσότερων ανωμαλιών σε ορισμένες ρυθμίσεις, χωρίς όμως να βελτιώνει απαραίτητα την ισορροπία μεταξύ \en{precision} και \en{recall}.

Συνολικά, η καλύτερη ομοσπονδιακή επίδοση ήταν το \en{F1-score} \(0.8245\), το οποίο προέκυψε με 4 \en{clients}, \en{FedAdam} και \(\alpha=100\). Η τιμή αυτή υπολείπεται της κεντρικής εκπαίδευσης κατά \(0.1249\) μονάδες \en{F1-score}. Η διαφορά είναι αναμενόμενη, επειδή στην ομοσπονδιακή εκπαίδευση κάθε \en{client} βλέπει μόνο το δικό του τμήμα των δεδομένων και ο διακομιστής συνδυάζει έμμεσα τις τοπικές ενημερώσεις. Αντίθετα, το κεντρικό μοντέλο βελτιστοποιείται απευθείας πάνω στο πλήρες σύνολο εκπαίδευσης.

Παρά την πτώση σε σχέση με την κεντρική εκπαίδευση, τα αποτελέσματα δείχνουν ότι η ομοσπονδιακή προσαρμογή με \en{LoRA} είναι λειτουργική για το συγκεκριμένο πρόβλημα. Οι περισσότερες ομοσπονδιακές εκτελέσεις συγκεντρώθηκαν σε στενό εύρος τιμών \en{F1-score} γύρω από το \(0.82\), χωρίς μεταφορά των αρχικών δεδομένων εκπαίδευσης στον κεντρικό διακομιστή. Το χαμηλότερο αποτέλεσμα, \(0.8011\), εμφανίστηκε στη ρύθμιση των 8 \en{clients} με \en{FedAvg} και \(\alpha=0.8\), όπου ο συνδυασμός μεγαλύτερου κατακερματισμού και πιο ανομοιογενούς κατανομής δυσκόλεψε περισσότερο τη συνολική σύγκλιση.

\section{Συμπέρασμα}

Το βασικό συμπέρασμα είναι ότι η κεντρική εκπαίδευση παραμένει σαφώς ισχυρότερη όταν η συγκέντρωση των δεδομένων είναι εφικτή. Η ομοσπονδιακή εκπαίδευση, ωστόσο, παρέχει μια πρακτική εναλλακτική όταν τα δεδομένα πρέπει να παραμείνουν κατανεμημένα, όπως συμβαίνει συχνά σε περιβάλλοντα δικτύων, παρόχων και συστημάτων όπου η ιδιωτικότητα, η ιδιοκτησία των δεδομένων ή το κόστος μεταφοράς αποτελούν περιοριστικούς παράγοντες.

Η μεταβολή του πλήθους των \en{clients} δεν οδήγησε σε μονοτονική βελτίωση της απόδοσης. Η καλύτερη τιμή \en{F1-score} αυξήθηκε ελαφρά από \(0.8201\) στους 2 \en{clients} σε \(0.8245\) στους 4 \en{clients}, αλλά μειώθηκε σε \(0.8222\) στους 8 \en{clients}. Επομένως, η αύξηση του πλήθους των \en{clients} δεν αρκεί από μόνη της για καλύτερη απόδοση. Αντίθετα, αυξάνει τον κατακερματισμό των τοπικών δεδομένων και μπορεί να κάνει τη συνάθροιση των ενημερώσεων πιο δύσκολη.

Η επίδραση της στρατηγικής συνάθροισης εξαρτήθηκε από τη ρύθμιση. Στους 2 \en{clients}, η καλύτερη επίδοση προήλθε από τη \en{FedAvg} με \(\alpha=100\). Στους 4 και 8 \en{clients}, η καλύτερη επίδοση προήλθε από τη \en{FedAdam} με \(\alpha=100\). Αυτό δείχνει ότι η προσαρμοστική ενημέρωση του διακομιστή γίνεται πιο χρήσιμη όταν αυξάνεται το πλήθος των τοπικών ενημερώσεων και η διαδικασία συνάθροισης γίνεται πιο σύνθετη.

Η παράμετρος \(\alpha\) επηρέασε επίσης τη συμπεριφορά των πειραμάτων, αλλά όχι με απόλυτα ενιαίο τρόπο. Οι καλύτερες τιμές \en{F1-score} και στις τρεις ομάδες \en{clients} προέκυψαν με \(\alpha=100\), δηλαδή με κατανομή που προσεγγίζει περισσότερο την \en{IID} περίπτωση. Ωστόσο, ορισμένες ανομοιογενείς ρυθμίσεις με \(\alpha=0.8\) πέτυχαν ανταγωνιστικές τιμές, ιδιαίτερα ως προς το \en{recall}. Αυτό δείχνει ότι η μη ομοιόμορφη κατανομή δεν οδηγεί αναγκαστικά σε αποτυχία, αλλά μπορεί να μεταβάλει την ισορροπία ανάμεσα στον εντοπισμό ανωμαλιών και στην αποφυγή ψευδώς θετικών προβλέψεων.

Η εργασία επιβεβαιώνει επίσης ότι η προσαρμογή \en{LoRA} μπορεί να ενσωματωθεί σε ομοσπονδιακό πλαίσιο για ανίχνευση ανωμαλιών σε αρχεία καταγραφής. Το μοντέλο δεν χρειάζεται να εκπαιδεύσει όλες τις παραμέτρους του \en{BERT}, αλλά μόνο τις πρόσθετες παραμέτρους προσαρμογής. Αυτό καθιστά την προσέγγιση πιο κατάλληλη για κατανεμημένα περιβάλλοντα, στα οποία το κόστος επικοινωνίας και οι υπολογιστικοί πόροι των \en{clients} αποτελούν σημαντικούς περιορισμούς.

Από πρακτική άποψη, τα αποτελέσματα δείχνουν έναν σαφή συμβιβασμό. Η κεντρική εκπαίδευση προσφέρει την καλύτερη απόδοση, αλλά απαιτεί συγκέντρωση των δεδομένων. Η ομοσπονδιακή εκπαίδευση μειώνει αυτή την απαίτηση, αλλά συνοδεύεται από χαμηλότερη επίδοση. Για εφαρμογές όπου η μέγιστη δυνατή ακρίβεια είναι ο μοναδικός στόχος, η κεντρική προσέγγιση παραμένει προτιμότερη. Για εφαρμογές όπου η ιδιωτικότητα, η κατανεμημένη ιδιοκτησία των δεδομένων και η περιορισμένη μεταφορά πληροφορίας είναι κρίσιμοι παράγοντες, η ομοσπονδιακή προσέγγιση αποτελεί ρεαλιστική λύση.

\section{Περιορισμοί και Μελλοντική Εργασία}

Η αξιολόγηση περιορίστηκε στο σύνολο δεδομένων \en{Loghub HDFS}, στο μοντέλο \en{bert-base-uncased} και σε ομοσπονδιακές διατάξεις με 2, 4 και 8 \en{clients}. Επίσης, εξετάστηκαν μόνο οι στρατηγικές \en{FedAvg} και \en{FedAdam}, δύο τιμές της παραμέτρου \(\alpha\) και ένα συγκεκριμένο σύνολο βασικών υπερπαραμέτρων για τα ομοσπονδιακά πειράματα. Μελλοντική εργασία μπορεί να επεκτείνει την αξιολόγηση σε περισσότερα σύνολα δεδομένων καταγραφών, διαφορετικά μοντέλα \en{transformer}, περισσότερες τιμές της παραμέτρου \(\alpha\), μεγαλύτερο πλήθος \en{clients} και πρόσθετες στρατηγικές ομοσπονδιακής βελτιστοποίησης, όπως \en{FedProx} ή προσωποποιημένες ομοσπονδιακές προσεγγίσεις.

Ένας ακόμη περιορισμός είναι ότι τα πειράματα βασίστηκαν σε σταθερό κοινό σύνολο επικύρωσης και ελέγχου. Αυτό επιτρέπει δίκαιη σύγκριση μεταξύ των εκτελέσεων, αλλά δεν εξετάζει περιπτώσεις όπου η κατανομή των δεδομένων αλλάζει με τον χρόνο ή όπου εμφανίζονται νέοι τύποι ανωμαλιών. Σε πραγματικά συστήματα καταγραφών, οι μορφές των συμβάντων και τα πρότυπα κανονικής λειτουργίας μπορούν να μεταβάλλονται. Επομένως, μελλοντική εργασία θα μπορούσε να εξετάσει συνεχή ή περιοδική ομοσπονδιακή ενημέρωση του μοντέλου, καθώς και μηχανισμούς προσαρμογής σε μεταβαλλόμενες κατανομές.

Επιπλέον, θα είχε αξία η ενσωμάτωση μηχανισμών ιδιωτικότητας και ασφάλειας, όπως η διαφορική ιδιωτικότητα και η ασφαλής συνάθροιση. Η παρούσα εργασία αποφεύγει τη μεταφορά των αρχικών δεδομένων στον διακομιστή, αλλά δεν μελετά επιπλέον προστασία απέναντι σε πιθανή διαρροή πληροφορίας μέσω των ενημερώσεων του μοντέλου. Τέτοιοι μηχανισμοί θα έκαναν το πειραματικό πλαίσιο πιο κοντά σε πραγματικές συνθήκες κατανεμημένης εκπαίδευσης.

Τέλος, θα ήταν χρήσιμη η λεπτομερέστερη ανάλυση του κόστους επικοινωνίας, του χρόνου εκπαίδευσης και της μεταβλητότητας των αποτελεσμάτων σε πολλαπλές επαναλήψεις. Η παρούσα εργασία εστιάζει κυρίως στις μετρικές ταξινόμησης, όμως σε ομοσπονδιακά συστήματα η απόδοση του μοντέλου πρέπει να αξιολογείται μαζί με το επικοινωνιακό και υπολογιστικό κόστος. Μια τέτοια ανάλυση θα επέτρεπε πιο ολοκληρωμένη σύγκριση ανάμεσα στην κεντρική και την ομοσπονδιακή προσέγγιση, καθώς και καλύτερη εκτίμηση του κατά πόσο η προτεινόμενη λύση μπορεί να εφαρμοστεί σε μεγαλύτερα και πιο ρεαλιστικά περιβάλλοντα.
